\chapter{Diferencat e fundme dhe simulimi i valëve}
\begin{qellime}
Studenti diskretizon derivatet hapësinore; ndërton skemën për ekuacionin e valës; zbaton kushte fillestare e kufitare; analizon kushtin CFL, dispersimin dhe reflektimin; interpreton valët bredhëse e të qëndrueshme.
\end{qellime}

\section{Nga PDE te rrjeti}
Në rrjetin $x_j=j\Delta x$, $t^n=n\Delta t$, fusha përfaqësohet nga $u_j^n\approx u(x_j,t^n)$. Për ekuacionin e valës
\begin{equation}
 \frac{\partial^2u}{\partial t^2}=c^2\frac{\partial^2u}{\partial x^2},
\end{equation}
diferencat qendrore japin
\begin{equation}
 u_j^{n+1}=2u_j^n-u_j^{n-1}+r^2(u_{j+1}^n-2u_j^n+u_{j-1}^n),
 \qquad r=\frac{c\Delta t}{\Delta x}.
\end{equation}
Skema është eksplicite dhe e rendit të dytë në kohë e hapësirë.

\section{Kushtet fillestare dhe kufitare}
Nevojiten zhvendosja $u(x,0)$ dhe shpejtësia $u_t(x,0)$. Hapi i parë ndërtohet me zhvillim Taylor. Skaji i fiksuar përdor Dirichlet, $u=0$; skaji i lirë përdor Neumann, $u_x=0$; një skaj thithës synon të minimizojë valën e reflektuar.

\begin{figure}[ht]
\centering\includegraphics[width=.82\textwidth]{12_vala_bredhese.pdf}
\caption{Evolucioni i një pulsi bredhës dhe përafrimi te kufiri reflektues.}
\end{figure}

\section{Stabiliteti CFL}
Analiza von Neumann zëvendëson $u_j^n$ me një mod Fourier $G^n e^{ikj\Delta x}$. Për skemën standarde në një përmasë merret kushti
\begin{equation}
 r=\frac{c\Delta t}{\Delta x}\le1.
\end{equation}
Ky kusht ka interpretim kauzal: domeni numerik i varësisë duhet të përmbajë domenin fizik. Shkelja e tij prodhon rritje eksponenciale të gabimeve.

\section{Dispersimi numerik}
Edhe kur skema është stabile, frekuenca numerike nuk përputhet saktësisht me $\omega=ck$. Për rrjetin hapësinor,
\begin{equation}
 k_{\mathrm{eff}}=\frac{2}{\Delta x}\sin\frac{k\Delta x}{2}.
\end{equation}
Modet me gjatësi vale pranë $2\Delta x$ përhapen gabim. Një rregull praktik kërkon shumë pika për gjatësi vale, jo vetëm dy.

\begin{figure}[ht]
\centering\includegraphics[width=.73\textwidth]{13_dispersimi_numerik.pdf}
\caption{Shpejtësia fazore numerike zvogëlohet për numra valorë të lartë.}
\end{figure}

\section{Valët e qëndrueshme}
Për skaje të fiksuara, zgjidhjet modale janë
\begin{equation}
 u_n(x,t)=A_n\sin\frac{n\pi x}{L}\cos(\omega_nt+\phi_n),\qquad
 \omega_n=\frac{n\pi c}{L}.
\end{equation}
Simulimi mund të nxjerrë spektrin përmes transformimit Fourier të sinjalit në një pikë. Krahasimi me $\omega_n$ verifikon rrjetin dhe kushtet kufitare.

\section{Reflektimi nga plazma}
Në një plazmë të ftohtë pa fushë magnetike,
\begin{equation}
 \varepsilon_r(\omega)=1-\frac{\omega_p^2}{\omega^2},\qquad
 \omega_p=\sqrt{\frac{n_e e^2}{m_e\varepsilon_0}}.
\end{equation}
Kur $\omega<\omega_p$, numri valor bëhet imagjinar dhe vala zbehet: plazma reflekton. Për $\omega>\omega_p$, vala përhapet me dispersim. Në model numerik, fusha elektromagnetike mund të lidhet me lëvizjen e elektroneve; ruajtja e energjisë dhe reflektimi teorik shërbejnë si kontrolle.

\begin{kujdes}
Një kufi i rrjetit mund të krijojë reflektim artificial që ngatërrohet me reflektimin fizik nga plazma. Së pari testohet skaji thithës në vakum, pastaj shtohet mjedisi dispersiv.
\end{kujdes}

\begin{lidhje}
Laboratori krijon valë bredhëse dhe të qëndrueshme, mat shpejtësinë, teston CFL dhe realizon një ndërfaqe vakum--plazmë. Studentët krahasojnë koeficientin numerik të reflektimit me parashikimin teorik.
\end{lidhje}

\section*{Pyetje kontrolli}
\begin{enumerate}
\item Çfarë dallimi ka stabiliteti nga dispersimi numerik?
\item Pse duhen më shumë se dy pika për gjatësi vale?
\item Si ndahet reflektimi fizik nga reflektimi artificial i kufirit?
\end{enumerate}
\section{Algoritmi me diferenca të fundme për valën}
\subsection{Zbatimi vektorial}
\begin{lstlisting}[language=Python,caption={Skema eksplicite e valës me skaje të fiksuara.}]
def vale_1d(u0, v0, c, dx, dt, n_hapa):
    r = c*dt/dx
    if r > 1.0:
        raise ValueError("Shkelet kushti CFL: c*dt/dx <= 1.")
    u0 = np.asarray(u0, float)
    u_prev = u0.copy()
    lap = u0[2:] - 2*u0[1:-1] + u0[:-2]
    u = u0.copy()
    u[1:-1] = u0[1:-1] + dt*v0[1:-1] + 0.5*r*r*lap
    u[[0, -1]] = 0.0
    historia = [u_prev.copy(), u.copy()]
    for _ in range(1, n_hapa):
        u_next = np.zeros_like(u)
        u_next[1:-1] = (2*u[1:-1] - u_prev[1:-1]
                          + r*r*(u[2:] - 2*u[1:-1] + u[:-2]))
        u_prev, u = u, u_next
        historia.append(u.copy())
    return np.asarray(historia)
\end{lstlisting}

\subsection{Kontrollet numerike}
Skema kontrollohet në katër nivele:
\begin{enumerate}
\item kufijtë plotësohen në çdo hap;
\item amplituda nuk rritet pa shkak për $r\le1$;
\item shpejtësia e pulsit afrohet me $c$ kur rafinohet rrjeta;
\item energjia diskrete mbetet afërsisht konstante për kufij reflektues.
\end{enumerate}
Një energji diskrete e dobishme është
\begin{equation}
 E^n=\frac12\sum_j\left(\frac{u_j^{n+1}-u_j^n}{\Delta t}\right)^2\Delta x
 +\frac{c^2}{2}\sum_j\left(\frac{u_{j+1}^n-u_j^n}{\Delta x}\right)^2\Delta x.
\end{equation}

\subsection{Matja e reflektimit}
Në dy sonda, një para dhe një pas ndërfaqes, regjistrohet $u(t)$. Amplituda incidente merret para mbërritjes së reflektimit, ndërsa amplituda e reflektuar në një dritare të vonuar. Koeficienti i amplitudës është $R_A=A_r/A_i$ dhe ai i energjisë $R_E\approx R_A^2$ vetëm kur impedancat janë të njëjta.

\begin{lstlisting}[language=Python]
def amplitude_dritare(sinjal, i0, i1):
    segment = sinjal[i0:i1]
    return 0.5*(segment.max() - segment.min())
\end{lstlisting}

\subsection{Reflektimi nga plazma}
Për një frekuencë të vetme, së pari llogaritet $\omega_p$ dhe parashikohet nëse vala duhet të përhapet. Simulimi nuk interpretohet para se të testohet vakumi dhe kufiri thithës. Në notebook, dendësia elektronike ndryshohet gradualisht dhe ndërtohet grafiku $R(\omega/\omega_p)$.

\begin{lidhje}
I njëjti funksion duhet të përdoret për valën bredhëse, valën e qëndrueshme dhe ndërfaqen e plazmës; ndryshojnë vetëm hyrjet dhe kushtet kufitare. Kjo e bën krahasimin fizik të kontrolluar.
\end{lidhje}

Diskretizimi dhe terminologjia e ekuacioneve me derivate të pjesshme mbështeten te tradita e analizës numerike në shqip \cite{hoxha,hanelli,naco}; pjesa e valëve dhe plazmës plotësohet me tekstet e fizikës llogaritëse \cite{giordano,landau}.

\subsection{Diskretizimi i ekuacionit të valës}
Për
\begin{equation}
 u_{tt}=c^2u_{xx},
\end{equation}
vendosim $x_j=j\Delta x$ dhe $t^n=n\Delta t$. Diferencat qendrore japin
\begin{equation}
 \frac{u_j^{n+1}-2u_j^n+u_j^{n-1}}{\Delta t^2}
 =c^2\frac{u_{j+1}^n-2u_j^n+u_{j-1}^n}{\Delta x^2}
 +\Oh(\Delta t^2+\Delta x^2).
\end{equation}
Pas zgjidhjes për shtresën e re,
\begin{equation}
 u_j^{n+1}=2u_j^n-u_j^{n-1}
 +\lambda^2(u_{j+1}^n-2u_j^n+u_{j-1}^n),
 \qquad \lambda=\frac{c\Delta t}{\Delta x}.
\end{equation}
Skema është eksplicite dhe e rendit të dytë në hapësirë e kohë. Për të nisur, nga $u_t(x,0)=v_0(x)$ merret
\begin{equation}
 u_j^1=u_j^0+\Delta t\,v_{0,j}
 +\frac{\lambda^2}{2}(u_{j+1}^0-2u_j^0+u_{j-1}^0).
\end{equation}

\subsection{Analiza von Neumann dhe kushti CFL}
Provojmë një mod Fourier $u_j^n=G^ne^{\mathrm{i}j\theta}$. Zëvendësimi në skemë jep
\begin{equation}
 G^2-2\left(1-2\lambda^2\sin^2\frac\theta2\right)G+1=0.
\end{equation}
Qëndrueshmëria kërkon që të dy rrënjët të kenë modul jo më të madh se një. Kjo ndodh kur
\begin{equation}
 \lambda=\frac{c\Delta t}{\Delta x}\le1.
\end{equation}
Interpretimi është kauzal: gjatë një hapi kohor, informacioni fizik përshkon largësinë $c\Delta t$, e cila nuk duhet të tejkalojë shtrirjen $\Delta x$ që skema përdor për komunikim ndërmjet nyjeve.

\subsection{Dispersimi numerik}
Duke vendosur $G=e^{-\mathrm{i}\omega\Delta t}$ merret lidhja e dispersimit
\begin{equation}
 \sin^2\left(\frac{\omega\Delta t}{2}\right)
 =\lambda^2\sin^2\left(\frac{k\Delta x}{2}\right).
\end{equation}
Në vazhdësi $\omega=ck$, por në rrjet raporti $\omega/k$ varet nga $k$. Modet me gjatësi vale të krahasueshme me $\Delta x$ përhapen me shpejtësi fazore të gabuar. Prandaj kushti CFL garanton qëndrueshmëri, jo saktësi. Një rregull praktik është të përdoren shumë pika për gjatësi vale dhe të kryhet studim rrjeti.

\subsection{Kushtet kufitare dhe energjia}
Kufiri i fiksuar përdor $u=0$ dhe reflekton me ndryshim shenje. Kufiri i lirë përdor $u_x=0$. Një kusht dalës i rendit të parë,
\begin{equation}
 u_t+cu_x=0,
\end{equation}
lejon kalimin e valës djathtas, por diskretizimi i tij mund të reflektojë pjesërisht. Energjia e vazhdueshme
\begin{equation}
 E(t)=\frac12\int\left(u_t^2+c^2u_x^2\right)\dd x
\end{equation}
është konstante për kufij pa fluks. Një energji diskrete e qëndrueshme është provë e fortë e implementimit.

\subsection{Valët dhe reflektimi nga plazma}
Zgjidhja e d'Alembert-it $u(x,t)=F(x-ct)+G(x+ct)$ ndan valët bredhëse djathtas e majtas. Në një interval me skaje të fiksuara, kushtet kufitare zgjedhin numra valorë diskretë $k_n=n\pi/L$ dhe frekuenca $\omega_n=ck_n$, duke formuar valë të qëndrueshme.

Në një plazmë të ftohtë pa fushë magnetike, permitiviteti relativ është
\begin{equation}
 \varepsilon_r(\omega)=1-\frac{\omega_p^2}{\omega^2},\qquad
 \omega_p=\sqrt{\frac{n_ee^2}{m_e\varepsilon_0}}.
\end{equation}
Për $\omega<\omega_p$, numri valor bëhet imagjinar dhe fusha shuhet brenda plazmës; vala reflektohet. Në modelin me përplasje, $\omega^2$ zëvendësohet në mënyrë efektive nga $\omega(\omega+\mathrm{i}\nu)$ dhe shfaqet absorbimi. Simulimi duhet të dallojë reflektimin fizik nga ai artificial i kufirit numerik.

\subsection{Përmbledhje dhe ushtrime}
Diskretizimi i valëve kërkon njëkohësisht konsistencë, CFL, rezolucion spektral dhe kushte kufitare fizikisht të përshtatshme.
\begin{enumerate}
\item Nxirrni formulën e shtresës së parë nga Taylor-i dhe shpejtësia fillestare.
\item Provoni kushtin CFL me analizën von Neumann.
\item Matni shpejtësinë fazore numerike për disa numra valorë.
\item Ndërtoni një valë të qëndrueshme dhe verifikoni frekuencat normale.
\item Simuloni kalimin nga vakumi në plazmë dhe ndani reflektimin fizik nga ai i kufirit të djathtë.
\end{enumerate}

\subsection{Koeficientët e reflektimit dhe transmetimit}
Kur një valë harmonike bie normalisht mbi një ndërfaqe, fusha në anën e parë shkruhet si shumë e valës rënëse dhe të reflektuar,
\begin{equation}
 E_1(z)=E_i e^{\mathrm{i}k_1z}+E_r e^{-\mathrm{i}k_1z},
\end{equation}
ndërsa në mjedisin e dytë $E_2(z)=E_t e^{\mathrm{i}k_2z}$. Vazhdimësia e komponentëve tangjencialë të fushave elektrike dhe magnetike në $z=0$ jep
\begin{equation}
 1+r=t,
 \qquad
 \frac{1-r}{Z_1}=\frac{t}{Z_2},
\end{equation}
ku $r=E_r/E_i$, $t=E_t/E_i$ dhe $Z_j$ është impedanca e mjedisit. Zgjidhja është
\begin{equation}
 r=\frac{Z_2-Z_1}{Z_2+Z_1},
 \qquad
 t=\frac{2Z_2}{Z_2+Z_1}.
\end{equation}
Koeficienti energjetik i reflektimit është $R=|r|^2$; transmetimi duhet të përfshijë raportin e flukseve dhe nuk është gjithmonë thjesht $|t|^2$.

Për plazmën e ftohtë pa përplasje, $\varepsilon_r=1-\omega_p^2/\omega^2$. Kur $\omega>\omega_p$, numri valor është real dhe vala transmetohet pjesërisht. Kur $\omega<\omega_p$, $k$ është imagjinar: fusha depërton vetëm me një gjatësi karakteristike dhe nuk transporton fuqi në thellësi në modelin ideal. Në plazmë me përplasje, permitiviteti bëhet kompleks; pjesa imagjinare përfaqëson humbjen dhe ndan absorbimin nga reflektimi.

\subsubsection{Ndërfaqja në rrjet dhe gabimi numerik}
Në diferencat e fundme, koeficienti material mund të ndryshojë ndërmjet nyjeve. Vendosja e ndërfaqes pikërisht në një nyje ose gjysmëhapi ndikon në rendin lokal. Një diskretizim konservativ i ekuacionit
\begin{equation}
 u_{tt}=\partial_x\left(c(x)^2u_x\right)
\end{equation}
përdor flukse në gjysmënyje,
\begin{equation}
 \left[\partial_x(c^2u_x)\right]_j
 \approx\frac{c_{j+1/2}^2(u_{j+1}-u_j)-c_{j-1/2}^2(u_j-u_{j-1})}{\Delta x^2}.
\end{equation}
Mesatarja e koeficientit në ndërfaqe duhet zgjedhur nga kushtet e vazhdimësisë; në probleme fluksi, mesatarja harmonike shpesh ruan më mirë fizikën se mesatarja aritmetike.

Për të matur reflektimin numerik, sinjali regjistrohet në një sondë para ndërfaqes. Simulimi referencë pa ndërfaqe jep pulsin rënës; diferenca izolon pulsin e reflektuar. Dritaret kohore duhet të ndajnë pulset dhe kufiri i majtë nuk duhet të reflektojë para përfundimit të matjes. Koeficienti mund të nxirret nga amplituda, energjia ose transformimi Fourier, por përkufizimi duhet të përputhet me formulën teorike.

\subsubsection{Absorbimi kufitar}
Një domen i fundëm duhet të imitojë hapësirën e pafundme. Kushti $u_t+cu_x=0$ është i saktë për valë njëdrejtimëshe në një përmasë dhe incidencë normale, por jo për të gjitha frekuencat dhe këndet. Një shtresë amortizuese shton termin $-\sigma(x)u_t$, me $\sigma$ që rritet gradualisht drejt kufirit. Rritja e menjëhershme e $\sigma$ krijon vetë një ndërfaqe dhe reflektim.

Shtresa e përputhur në mënyrë të përsosur (PML) projektohet që impedanca të mbetet e përputhur në hyrje dhe vala të shuhet brenda shtresës. Implementimi i saj është më i ndërlikuar, por shërben si shembull se kushti kufitar është pjesë e modelit numerik. Në laborator mund të krahasohen kufiri i fiksuar, kushti dalës i rendit të parë dhe shtresa amortizuese duke matur energjinë e kthyer.

\subsubsection{Bilanci diskret i energjisë}
Për skemën qendrore, shpejtësia diskrete mund të vendoset në gjysmëhapa. Një energji e përshtatshme është
\begin{equation}
 E^{n+1/2}=\frac{\Delta x}{2}\sum_j
 \left(\frac{u_j^{n+1}-u_j^n}{\Delta t}\right)^2
 +\frac{c^2\Delta x}{2}\sum_j
 \left(\frac{u_{j+1}^{n+1}-u_j^{n+1}}{\Delta x}\right)
 \left(\frac{u_{j+1}^{n}-u_j^{n}}{\Delta x}\right).
\end{equation}
Për kufij të përshtatshëm dhe CFL të plotësuar, kjo madhësi mbetet e kufizuar. Monitorimi i saj zbulon paqëndrueshmërinë shumë më herët se inspektimi vizual i animacionit.

